% ============================================================
%  Elegant 数学教学文档 · 共享导言区
%  用法：在 \documentclass{elegantbook / elegantnote} 之后
%        % ============================================================
%  Elegant 数学教学文档 · 共享导言区
%  用法：在 \documentclass{elegantbook / elegantnote} 之后
%        % ============================================================
%  Elegant 数学教学文档 · 共享导言区
%  用法：在 \documentclass{elegantbook / elegantnote} 之后
%        % ============================================================
%  Elegant 数学教学文档 · 共享导言区
%  用法：在 \documentclass{elegantbook / elegantnote} 之后
%        \input{../common/preamble}
%  依赖：TeX Live 2026 + XeLaTeX
% ============================================================

% ---------- 等宽字体 ----------
% lmmono10 缺少 λ、≥ 等符号，改用覆盖更全的 DejaVu Sans Mono
\setmonofont{DejaVu Sans Mono}

% ---------- 数学 ----------
\usepackage{mathtools}   % amsmath 扩展（\substack、\dcases、\prescript 等）
\usepackage{bm}          % 粗体数学符号 \bm
\usepackage{cases}       % 分段函数 / 方程组 \begin{cases}
\usepackage{cancel}      % 划去 \cancel、\cancelto
\usepackage{siunitx}     % 单位与数字 \num \SI
\sisetup{
  detect-all,
  per-mode = symbol,
  range-units = single,
}

% ---------- 图形 ----------
\usepackage{tikz}
\usetikzlibrary{arrows.meta, calc, decorations.pathreplacing, positioning}
\usepackage{pgfplots}
\pgfplotsset{compat=1.18}

% ---------- 统一颜色名 ----------
% elegantbook 定义了 main / second / third；elegantnote 只有 ecolor。
% 用 \providecolor 补齐，使共享宏在两个模板下都能用。
\providecolor{main}{RGB}{0,120,2}
\providecolor{second}{RGB}{230,90,7}
\providecolor{third}{RGB}{0,160,152}

% ---------- 版式 ----------
\usepackage{fontawesome5}   % 图标（如无需要可删除本行）

% ---------- 数学字体补丁 ----------
% elegantbook/elegantnote 的 math=newtx 已加载 newtxmath；
% 这里补上常用但 newtx 未覆盖的符号字体。
\usepackage{mathrsfs}       % \mathscr

%  依赖：TeX Live 2026 + XeLaTeX
% ============================================================

% ---------- 等宽字体 ----------
% lmmono10 缺少 λ、≥ 等符号，改用覆盖更全的 DejaVu Sans Mono
\setmonofont{DejaVu Sans Mono}

% ---------- 数学 ----------
\usepackage{mathtools}   % amsmath 扩展（\substack、\dcases、\prescript 等）
\usepackage{bm}          % 粗体数学符号 \bm
\usepackage{cases}       % 分段函数 / 方程组 \begin{cases}
\usepackage{cancel}      % 划去 \cancel、\cancelto
\usepackage{siunitx}     % 单位与数字 \num \SI
\sisetup{
  detect-all,
  per-mode = symbol,
  range-units = single,
}

% ---------- 图形 ----------
\usepackage{tikz}
\usetikzlibrary{arrows.meta, calc, decorations.pathreplacing, positioning}
\usepackage{pgfplots}
\pgfplotsset{compat=1.18}

% ---------- 统一颜色名 ----------
% elegantbook 定义了 main / second / third；elegantnote 只有 ecolor。
% 用 \providecolor 补齐，使共享宏在两个模板下都能用。
\providecolor{main}{RGB}{0,120,2}
\providecolor{second}{RGB}{230,90,7}
\providecolor{third}{RGB}{0,160,152}

% ---------- 版式 ----------
\usepackage{fontawesome5}   % 图标（如无需要可删除本行）

% ---------- 数学字体补丁 ----------
% elegantbook/elegantnote 的 math=newtx 已加载 newtxmath；
% 这里补上常用但 newtx 未覆盖的符号字体。
\usepackage{mathrsfs}       % \mathscr

%  依赖：TeX Live 2026 + XeLaTeX
% ============================================================

% ---------- 等宽字体 ----------
% lmmono10 缺少 λ、≥ 等符号，改用覆盖更全的 DejaVu Sans Mono
\setmonofont{DejaVu Sans Mono}

% ---------- 数学 ----------
\usepackage{mathtools}   % amsmath 扩展（\substack、\dcases、\prescript 等）
\usepackage{bm}          % 粗体数学符号 \bm
\usepackage{cases}       % 分段函数 / 方程组 \begin{cases}
\usepackage{cancel}      % 划去 \cancel、\cancelto
\usepackage{siunitx}     % 单位与数字 \num \SI
\sisetup{
  detect-all,
  per-mode = symbol,
  range-units = single,
}

% ---------- 图形 ----------
\usepackage{tikz}
\usetikzlibrary{arrows.meta, calc, decorations.pathreplacing, positioning}
\usepackage{pgfplots}
\pgfplotsset{compat=1.18}

% ---------- 统一颜色名 ----------
% elegantbook 定义了 main / second / third；elegantnote 只有 ecolor。
% 用 \providecolor 补齐，使共享宏在两个模板下都能用。
\providecolor{main}{RGB}{0,120,2}
\providecolor{second}{RGB}{230,90,7}
\providecolor{third}{RGB}{0,160,152}

% ---------- 版式 ----------
\usepackage{fontawesome5}   % 图标（如无需要可删除本行）

% ---------- 数学字体补丁 ----------
% elegantbook/elegantnote 的 math=newtx 已加载 newtxmath；
% 这里补上常用但 newtx 未覆盖的符号字体。
\usepackage{mathrsfs}       % \mathscr

%  依赖：TeX Live 2026 + XeLaTeX
% ============================================================

% ---------- 等宽字体 ----------
% lmmono10 缺少 λ、≥ 等符号，改用覆盖更全的 DejaVu Sans Mono
\setmonofont{DejaVu Sans Mono}

% ---------- 数学 ----------
\usepackage{mathtools}   % amsmath 扩展（\substack、\dcases、\prescript 等）
\usepackage{bm}          % 粗体数学符号 \bm
\usepackage{cases}       % 分段函数 / 方程组 \begin{cases}
\usepackage{cancel}      % 划去 \cancel、\cancelto
\usepackage{siunitx}     % 单位与数字 \num \SI
\sisetup{
  detect-all,
  per-mode = symbol,
  range-units = single,
}

% ---------- 图形 ----------
\usepackage{tikz}
\usetikzlibrary{arrows.meta, calc, decorations.pathreplacing, positioning}
\usepackage{pgfplots}
\pgfplotsset{compat=1.18}

% ---------- 统一颜色名 ----------
% elegantbook 定义了 main / second / third；elegantnote 只有 ecolor。
% 用 \providecolor 补齐，使共享宏在两个模板下都能用。
\providecolor{main}{RGB}{0,120,2}
\providecolor{second}{RGB}{230,90,7}
\providecolor{third}{RGB}{0,160,152}

% ---------- 版式 ----------
\usepackage{fontawesome5}   % 图标（如无需要可删除本行）

% ---------- 数学字体补丁 ----------
% elegantbook/elegantnote 的 math=newtx 已加载 newtxmath；
% 这里补上常用但 newtx 未覆盖的符号字体。
\usepackage{mathrsfs}       % \mathscr
